% Options for packages loaded elsewhere
\PassOptionsToPackage{unicode}{hyperref}
\PassOptionsToPackage{hyphens}{url}
\PassOptionsToPackage{dvipsnames,svgnames,x11names}{xcolor}
%
\documentclass[
  11pt,
]{article}
\usepackage{amsmath,amssymb}
\usepackage[]{mathptmx}
\usepackage{iftex}
\ifPDFTeX
  \usepackage[T1]{fontenc}
  \usepackage[utf8]{inputenc}
  \usepackage{textcomp} % provide euro and other symbols
\else % if luatex or xetex
  \usepackage{unicode-math}
  \defaultfontfeatures{Scale=MatchLowercase}
  \defaultfontfeatures[\rmfamily]{Ligatures=TeX,Scale=1}
\fi
% Use upquote if available, for straight quotes in verbatim environments
\IfFileExists{upquote.sty}{\usepackage{upquote}}{}
\IfFileExists{microtype.sty}{% use microtype if available
  \usepackage[]{microtype}
  \UseMicrotypeSet[protrusion]{basicmath} % disable protrusion for tt fonts
}{}
\makeatletter
\@ifundefined{KOMAClassName}{% if non-KOMA class
  \IfFileExists{parskip.sty}{%
    \usepackage{parskip}
  }{% else
    \setlength{\parindent}{0pt}
    \setlength{\parskip}{6pt plus 2pt minus 1pt}}
}{% if KOMA class
  \KOMAoptions{parskip=half}}
\makeatother
\usepackage{xcolor}
\usepackage[margin=1in]{geometry}
\setlength{\emergencystretch}{3em} % prevent overfull lines
\providecommand{\tightlist}{%
  \setlength{\itemsep}{0pt}\setlength{\parskip}{0pt}}
\setcounter{secnumdepth}{-\maxdimen} % remove section numbering
\ifLuaTeX
  \usepackage{selnolig}  % disable illegal ligatures
\fi
\IfFileExists{bookmark.sty}{\usepackage{bookmark}}{\usepackage{hyperref}}
\IfFileExists{xurl.sty}{\usepackage{xurl}}{} % add URL line breaks if available
\urlstyle{same} % disable monospaced font for URLs
\hypersetup{
  colorlinks=true,
  linkcolor={blue},
  filecolor={Maroon},
  citecolor={Blue},
  urlcolor={blue},
  pdfcreator={LaTeX via pandoc}}

\author{}
\date{}

\begin{document}

\hypertarget{the-organisation-as-the-executable-unit-towards-an-executable-organisational-operating-model-for-ai-native-software-engineering}{%
\section{The Organisation as the Executable Unit: Towards an Executable
Organisational Operating Model for AI-Native Software
Engineering}\label{the-organisation-as-the-executable-unit-towards-an-executable-organisational-operating-model-for-ai-native-software-engineering}}

\textbf{Sudha Manimaran} \emph{Independent Researcher}

Version 1.0 · August 2026

\hypertarget{abstract}{%
\subsection{Abstract}\label{abstract}}

Software engineering has historically depended on organisational
structures, shared objectives, delegated authority, reviewed decisions,
and institutional knowledge that outlives any individual contributor, to
sustain systems across changes in personnel and technology. Contemporary
artificial intelligence systems increasingly perform substantial
engineering work, including code generation, architectural proposal, and
review, yet operate largely outside these organisational structures:
without delegated authority, without persistent organisational
responsibility, and without accountability that survives the interaction
in which a decision was made. The paper argues that this absence of
organisational structure, not insufficient model capability, is the
primary obstacle to safely extending AI autonomy across the software
engineering lifecycle. It proposes an executable organisational
operating model in which the organisation itself is the unit that
persists, not the individual participant or the AI agent. Objectives
justify capabilities. Capabilities expose governed services. Services
are delivered through accountable roles. Roles are fulfilled by
interchangeable participants. Those participants may be human,
artificial, or external. They produce artefacts, evidence, and
organisational knowledge. The paper distinguishes this approach from
existing agent-orchestration and workflow-coordination systems, which
coordinate interactions and activities but do not execute an
organisation with explicit authority, accountability, and institutional
memory. It discusses the conditions under which this becomes necessary
as the number of autonomous participants in an engineering organisation
grows, the limitations of governance as a guarantee of process rather
than judgement, and the extension of the model beyond software
engineering to other knowledge-intensive domains.

\hypertarget{introduction}{%
\subsection{1. Introduction}\label{introduction}}

Software engineering has never been an individual activity. Enduring
software systems are produced by organisations, whose capability arises
not from employing highly capable individuals but from coordinating
specialised responsibilities through governance, accountability, and
accumulated knowledge. Authority is granted rather than assumed.
Responsibilities attach to organisational roles rather than to the
individuals who currently occupy them. Decisions are reviewed before
they become organisational commitments. Knowledge is preserved
independently of the individuals who created it. These principles were
not designed in the abstract; they exist because they let software
systems outlive individual engineers, individual projects, and
successive generations of technology.

Existing software methodologies were designed around a human workforce
and its characteristics. Artificial intelligence changes that premise.
Current AI systems automate substantial engineering activity, producing
software, documentation, and architectural proposals of real quality,
without participating in the organisational structures that make
engineering sustainable. They hold no delegated authority, no
organisational responsibility, and no accountability that outlives the
interaction in which a decision was made. Once that interaction
concludes, the reasoning behind the decision typically disappears with
it rather than becoming part of an organisation's institutional
knowledge.

The paper argues that this is the central obstacle to extending AI
autonomy safely across the engineering lifecycle, and that it is an
organisational problem rather than a capability problem. The industry's
response to date has concentrated almost entirely on increasing what
individual AI systems can do: coding assistants evolving into autonomous
engineering agents, and a rapidly growing set of orchestration
frameworks coordinating increasingly specialised systems (Chen et al.,
2024). That focus addresses the participant. It does not address the
organisation the participant operates inside, a gap that recent work on
accountability in AI-powered software development has begun to document
from a different angle, examining how responsibility is currently
assigned, largely through vendor terms of service rather than
organisational structure (Treude, 2026).

\hypertarget{why-coordination-mechanisms-are-not-organisations}{%
\subsection{2. Why Coordination Mechanisms Are Not
Organisations}\label{why-coordination-mechanisms-are-not-organisations}}

Prompts coordinate individual interactions. Workflows coordinate
sequences of activity. Agent-orchestration frameworks coordinate
computational participants. None of these mechanisms executes an
organisation in the sense the term is used here: none represents
explicit, persistent authority, accountability, or institutional memory
that survives the interaction, the workflow run, or the individual agent
that happens to be participating at a given moment. This is consistent
with classical coordination theory, which treats coordination as the
management of dependencies among activities (Malone \& Crowston, 1994);
coordinating activities is a narrower problem than governing the
organisation those activities occur within, and the two should not be
conflated.

This distinction matters because the failure modes of ungoverned AI
participation are organisational failure modes, not capability failure
modes. A decision made independently by an AI agent, however
well-reasoned, carries no accountable chain behind it if nothing
external to that agent recorded who or what approved it, on what
evidence, and under what authority. Without that record, an organisation
cannot confirm that a given failure will not recur, cannot reconstruct
why a decision was made once the participant that made it has been
replaced, and cannot distinguish a well-justified engineering choice
from an unreviewed one after the fact.

Human engineering organisations solve this problem through mechanisms
refined over decades: authority that can be withheld as well as granted,
work checked by someone other than whoever performed it, and a record
that persists independently of any single individual (Finkelstein,
2009). This paper takes the position that these mechanisms should apply
to artificial participants for the same reason they apply to human ones,
not because artificial intelligence is untrustworthy, but because no
participant, human or artificial, should be able to unilaterally certify
its own work within an organisation that depends on that certification
being reliable. Recent work reframing autonomous agents as
organisational actors subject to agency theory and decision-rights
frameworks arrives at a related conclusion from corporate-governance
theory rather than software engineering practice (Chinnaraju, 2026);
this paper takes software engineering practice, rather than governance
theory in general, as its starting point.

\hypertarget{an-executable-organisational-operating-model}{%
\subsection{3. An Executable Organisational Operating
Model}\label{an-executable-organisational-operating-model}}

The paper proposes modelling the organisation as a persistent structure
whose identity does not depend on which participants currently occupy
it. Organisations derive continuity not from the people who currently
work within them but from structures that persist despite continual
changes in personnel: objectives remain, responsibilities remain,
standards remain, governance continues, and organisational knowledge
accumulates, even as individual participants change. The same principle,
this paper argues, should hold when some or all of an organisation's
participants are artificial.

The term executable requires precise definition, because executable
process already has an established meaning. That established meaning,
the sense in which BPMN and workflow engines are executable, differs
from what is intended here and could otherwise be mistaken for it. A
workflow engine executes a process: it interprets a description of
activities and control flow, and enforces the order in which tasks
occur. It does not represent, and therefore cannot enforce, the
organisational constructs that determine whether a given transition is
legitimate independent of sequencing, who holds authority to perform an
action, who holds the separate authority to approve it, what evidence a
state change requires before it is accepted, or how the outcome of that
check feeds back into the standards future transitions are evaluated
against. This paper defines an organisational model as executable if its
organisational constructs, authority constraints, accountability
relationships, service obligations, evidence requirements, and
governance rules, are represented in a form that a computational system
interprets and enforces during operation, rather than merely documented
for human reference and applied at human discretion. On this definition,
executability is a property of an organisation's authority and
accountability structure, not of its process sequencing; a system can
execute a workflow without executing an organisation in this sense,
which is the gap Section 2 argues current AI tooling has not closed.

The model represents the organisation as a chain of derived entities.
Objectives justify the existence of capabilities: an organisation
requires certain abilities because it has certain purposes it must
fulfil, not the reverse. Capabilities expose governed services, stable
interfaces through which one part of the organisation provides value to
another without exposing how that value is internally produced. Services
are delivered through accountable roles, which establish sustained
organisational responsibility independent of who or what currently
fulfils them. Roles are fulfilled by participants, human engineers, AI
agents, or external systems, who perform the activities required to
realise those responsibilities and produce artefacts, evidence, and,
ultimately, organisational knowledge.

This sequence is significant because it separates concepts that are
traditionally conflated in software engineering practice. Objectives are
not capabilities. Capabilities are not departments. Services are not
activities. Roles are not participants. Participants are not
technologies. Each of these is a distinct organisational abstraction
with a different rate of change: objectives evolve slowly, capabilities
more slowly still, while participants and their underlying technologies
may change continuously. Separating these layers allows an organisation
to evolve technologically without repeatedly redefining itself
organisationally, and allows successive generations of artificial
intelligence to be adopted as they mature without requiring the
organisation's governance, accountability, or institutional knowledge to
be rebuilt each time.

The chain described above is not static; it closes on itself through a
learning loop. Every service interaction generates evidence of its own
effectiveness. Evidence accumulates into organisational knowledge.
Organisational knowledge refines the capability definitions, standards,
and governance rules through which future services are delivered, so
that improved capabilities produce higher-quality services in turn:
capability generates service, service generates evidence, evidence
generates knowledge, and knowledge strengthens capability. This closed
loop is what makes the model adaptive rather than merely stable.

The loop also makes explicit a distinction that is otherwise easy to
elide: organisational learning is not the same as participant learning.
Replacing an AI participant with a more capable underlying model
improves that participant; it does not, by itself, improve the
organisation. Organisational learning occurs only when evidence changes
the operating model itself, a revised review standard, a tightened
governance rule, an updated capability definition, a change that
persists independently of which participant is subsequently assigned to
the corresponding role. This distinction matters precisely because the
paper has argued that the operating model, not the participant, is the
unit that persists; a learning mechanism confined to participants would
leave that unit unchanged no matter how much individual performance
improved.

\hypertarget{authority-evidence-and-the-limits-of-self-certification}{%
\subsection{4. Authority, Evidence, and the Limits of
Self-Certification}\label{authority-evidence-and-the-limits-of-self-certification}}

Within this model, two mechanisms do the load-bearing work of
accountability. First, the authority to perform a piece of work and the
authority to approve it are distinct grants; no participant holds both
for the same piece of work, regardless of whether that participant is
human, artificial, or an external system. Second, completion is a
proposal rather than a fact: a defined requirement, a passed test, a
completed review, or a recorded decision, must be satisfied and
independently checked before work is considered done. Neither mechanism
is novel in isolation; both are long-standing practice in mature human
engineering organisations. What the paper argues is novel is applying
them to artificial participants by default and as a structural property
of the organisation, rather than as a policy that depends on a human
remembering to enforce it, which is the condition under which most
current AI tooling operates.

It is important to state plainly what this buys and what it does not.
Structural governance of this kind proves that a process was followed.
It does not prove that the underlying judgement was good. An AI
participant approving another AI participant's work can clear every
gate, evidenced and traceable, and still be wrong. The claim made for
this model is auditability, not correctness: it ensures that engineering
decisions remain explainable, reviewable, and attributable, not that
they are always right. The paper treats this as an important and
underexamined distinction for any system that governs autonomous
participants.

\hypertarget{scale-as-the-forcing-function}{%
\subsection{5. Scale as the Forcing
Function}\label{scale-as-the-forcing-function}}

A handful of AI participants can be coordinated adequately through
prompts, workflows, and direct human supervision. The paper argues this
coordination model does not scale, and that the point at which it breaks
is organisationally, not technically, determined. As autonomy increases,
a typical engineering organisation is likely to manage not one or two AI
systems but many, specialised across requirements, architecture,
implementation, verification, security, compliance, and operations.
Coordinating that many autonomous participants simultaneously is not
solved by improving any individual participant's competence; it requires
an organisational abstraction for authority, accountability, and
institutional knowledge that operates at that scale. This is consistent
with a broader pattern in enterprise computing, in which new
coordination infrastructure becomes necessary not when a new technology
becomes possible but when an existing coordination mechanism breaks
under scale.

\hypertarget{generalisation-beyond-software-engineering}{%
\subsection{6. Generalisation Beyond Software
Engineering}\label{generalisation-beyond-software-engineering}}

The model as described, objectives, capabilities, governed services,
accountable roles, evidence, and organisational learning, is not
expressed in terms specific to software. Software engineering is, this
paper suggests, the first domain in which this problem becomes visible,
because it is the first knowledge-intensive discipline in which
artificial participants have become capable enough to expose the absence
of an organisational model around them. Whether the theory holds with
the same rigour in other knowledge-intensive domains, healthcare,
financial services, legal services, scientific research, has not been
established, and it is not presented here as a settled claim. It is
treated as a plausible extension worth testing empirically in future
work, not as a conclusion this paper argues for.

\hypertarget{limitations}{%
\subsection{7. Limitations}\label{limitations}}

The paper notes three limitations to the position argued here. First, as
stated in Section 4, structural governance provides auditability, not a
guarantee of correctness; an organisation adopting this model still
depends on the underlying competence of its participants and the quality
of its evidence standards. Second, the model as described requires an
organisation's capabilities, services, and dependency structure to be
explicitly authored before engineering work can run through it; this is
nontrivial setup work whose cost has not been characterised empirically
here. Third, the claims in Sections 5 and 6 concerning scale and
cross-domain generality are argued from first principles and from the
internal logic of the model; they have not yet been validated against
longitudinal data from organisations operating at the scale described,
and that validation is treated as necessary future work.

\hypertarget{conclusion}{%
\subsection{8. Conclusion}\label{conclusion}}

This paper has argued that the central obstacle to extending AI autonomy
safely across software engineering is organisational rather than
technical: artificial participants increasingly perform substantial
engineering work without operating inside organisational structures that
grant delegated authority, enforce independent review, or preserve
institutional knowledge beyond the interaction in which a decision was
made. It proposed an executable organisational operating model in which
the organisation, rather than any individual participant, is the unit
that persists, and in which authority, evidence, and accountability
apply to artificial participants by default rather than as optional
policy. It distinguished this from prompt-based and workflow-based
coordination mechanisms, which do not execute an organisation in this
sense, and outlined the conditions, principally scale, under which this
distinction is expected to become organisationally necessary rather than
merely desirable.

\hypertarget{references}{%
\subsection{References}\label{references}}

Chen, S., Liu, Y., Han, W., Zhang, W., \& Liu, T. (2024). A survey on
LLM-based multi-agent system: Recent advances and new frontiers in
application. \emph{arXiv preprint arXiv:2412.17481}.

Chinnaraju, A. (2026). When AI agents act: Governance, accountability,
and strategic risk in autonomous organizations. \emph{International
Journal of Research and Scientific Innovation}, 12(12).

Finkelstein, A. (2009). \emph{Software engineering governance: A
briefing}. Tutorial presented at the 31st International Conference on
Software Engineering (ICSE 2009).

Malone, T. W., \& Crowston, K. (1994). The interdisciplinary study of
coordination. \emph{ACM Computing Surveys}, 26(1), 87--119.
https://doi.org/10.1145/174666.174668

Treude, C. (2026). Accountable agents in software engineering: An
analysis of terms of service and a research roadmap. \emph{Proceedings
of the 3rd ACM International Conference on AI-Powered Software (AIware
2026)}. arXiv:2605.04532.

\end{document}
